\documentclass[12pt,a4paper]{article}

% === PACKAGES GÉNÉRAUX ===
\usepackage[utf8]{inputenc}      % Encodage UTF-8
\usepackage[T1]{fontenc}         % Codage de la police
\usepackage[english]{babel}      % Langue du document
\usepackage{lmodern}             % Police vectorielle
\usepackage{amsmath,amssymb}     % Maths propres
\usepackage{graphicx}            % Pour les images
\usepackage{booktabs}            % Pour de belles tables
\usepackage{setspace}            % Pour gérer l'interligne
\usepackage{hyperref}            % Liens cliquables dans le PDF
\usepackage{xcolor}              % Couleurs pour les liens
\usepackage{geometry}            % Contrôle des marges
\usepackage{enumitem}            % Listes personnalisées

% === MISE EN PAGE ===
\geometry{
    a4paper,
    left=2.5cm,
    right=2.5cm,
    top=2.5cm,
    bottom=2.5cm,
    headsep=10pt
}

\setstretch{1.2} % Interligne légèrement augmenté

% === LIENS ET COULEURS ===
\hypersetup{
    colorlinks=true,
    linkcolor=blue!50!black,
    urlcolor=blue!60!black,
    citecolor=blue!60!black,
    pdftitle={Winter Bootcamp},
    pdfauthor={Nathan Sauldubois}
}

% === TITRE ET MÉTADONNÉES ===
\title{\textbf{Winter Bootcamp Syllabus}\\[0.5em]
\large Mathematical Foundations for Finance and Risk Engineering}
\author{Nathan Sauldubois}
\date{The bootcamp will take place during the weeks of January 5 and January 12, running \textbf{from Monday, January 5 through Friday, January 16}.
\\
Sessions will be held on \textbf{Monday, Wednesday, and Friday} from \textbf{8:30 a.m. to 12:15 p.m.} .}

\begin{document}

\maketitle


The purpose of this Bootcamp is to provide you with the mathematical tools necessary to handle most of the quantitative concepts you will encounter during the program. 
The topics covered in this Bootcamp are: \textbf{Matrices, Differential Calculus, Optimization, and Probability.}


Except for the first session, which will be entirely dedicated to matrices, each lecture will consist of \textbf{1 hour and 30 minutes} of exercise review (based on assignments given from one class to the next), followed by a \textbf{15-minute break}, and then a \textbf{2-hour lecture}. 

All slides and exercises will be made available online. 
\textbf{Standard exercises} will be corrected during class, 
and I will upload \textbf{selected elements of correction} for the advanced exercises.

There will be \textbf{no final exam} for this refresher course, nor any \textbf{graded homework assignments}. 
However, it is \textbf{highly recommended} that you complete the \emph{standard exercises} between sessions, 
as they will help you quickly assimilate the concepts studied during the lectures.

\medskip 
If you have any questions, feel free to contact me by email at: \textbf{ns6982@nyu.edu}.
\medskip 



\section{Matrices (Lecture 1)}
The first part of the Bootcamp focuses on the fundamental operations involving matrices and their applications to finance. 
We will review how to manipulate matrices, solve linear systems, diagonalize matrices and see how it can be used, either to understand the role of covariance matrices in portfolio theory, solve an ODE linked with interest rate models.

\medskip 

\textbf{Main concepts:}
\begin{itemize}[itemsep=-4pt, topsep=2pt]
    \item Definition of matrices and operations on matrices
    \item Transpose (and symetric matrices), determinant, and inverse.
    \item Linear systems: solving $A x = b$ and interpretation.
    \item Eigenvalues, eigenvectors, and diagonalization (intuitive approach) including spectral theorem.
    \item Covariance and correlation matrices with applications.
\end{itemize}

\textbf{Applications:}
\begin{itemize}[itemsep=-4pt, topsep=2pt]
    \item Covariance and correlation matrices with applications, introduction to Portfolio variance: $\text{Var}(R_p) = w^{\top}\Sigma w$ and minimum variance portfolios: solving $\Sigma w = \lambda \mathbf{1}$.
    \item Risk factor decomposition and Principal Component Analysis (PCA).
    \item Matrix exponential $e^{A t}$ and applications to dynamic systems (interest rate model, Vasicek model).
\end{itemize}

\textbf{Suggested references to go further}
\begin{itemize}[itemsep=-4pt, topsep=2pt]
    \item Gilbert Strang, \textit{Introduction to Linear Algebra}.
    \item David C. Lay, \textit{Linear Algebra and Its Applications}.
    \item J. Douglas Carroll and Paul E. Green \textit{Mathematical Tools for Applied Multivariate Analysis}.
\end{itemize}

\section{Differential Calculus (Lecture 2)}
This section reviews single and multivariate differential calculus and emphasizes how these concepts appear in optimization, pricing models, and sensitivity analysis.

\textbf{Main concepts:}
\begin{itemize}[itemsep=-4pt, topsep=2pt]
    \item Derivatives and partial derivatives.
    \item Integration and derivation of integrals
    \item Gradient and Hessian matrices.
    \item Taylor expansion and approximation.
    \item Differentiability of vector-valued functions.
\end{itemize}

\textbf{Applications:}
\begin{itemize}[itemsep=-4pt, topsep=2pt]
    \item Heat equation, linear-quadratic HJB, $1$ dimensional transport equation and their solutions.
    \item Option pricing and greeks.
    \item Identifying convexity via the Hessian.
\end{itemize}

\textbf{Suggested references:}
\begin{itemize}[itemsep=-4pt, topsep=2pt]
    \item Carl P. Simon and Lawrence Blume, \textit{Mathematics for Economists}.
    \item Sydsaeter and Hammond, \textit{Essential Mathematics for Economic Analysis}.
    \item J. Douglas Carroll and Paul E. Green \textit{Mathematical Tools for Applied Multivariate Analysis}.
\end{itemize}


\section{Optimization (Lecture 2)}
Optimization plays a central role in financial decision-making, from portfolio selection to risk minimization. 
This part covers both unconstrained and constrained optimization problems and introduces Lagrange multipliers.

\textbf{Main concepts:}
\begin{itemize}[itemsep=-4pt, topsep=2pt]
    \item First- and second-order conditions for optimality.
    \item Convex and concave functions.
    \item Lagrange multipliers and constraint qualification.
\end{itemize}

\textbf{Applications:}
\begin{itemize}[itemsep=-4pt, topsep=2pt]
    \item Mean-variance optimization (Markowitz portfolio theory).
    \item Budget constraints and shadow prices.
    \item Calibration problems and least-squares estimation.
    \item Linear Regression
\end{itemize}

\textbf{Suggested references:}
\begin{itemize}[itemsep=-4pt, topsep=2pt]
    \item Boyd and Vandenberghe, \textit{Convex Optimization}.
    \item J. Douglas Carroll and Paul E. Green \textit{Mathematical Tools for Applied Multivariate Analysis}.
\end{itemize}


\section{Probability}

This final module of the Bootcamp is devoted to the fundamental concepts of probability theory and their direct applications. The goal is to understand how uncertainty is modeled mathematically and how probabilistic tools support portfolio analysis, pricing, and risk measurement.

\subsection{Basics (Lecture 3)}
We begin with the foundations of probability:
\begin{itemize}[itemsep=-4pt, topsep=2pt]
    \item Definition of a probability space and events.
    \item Probability rules, independence, and conditional probability.
    \item Law of total probability and Bayes' theorem.
\end{itemize}
Applications include modeling discrete events such as default risk, and interpreting joint and conditional probabilities in financial contexts.

\subsection{Random Variables (Lectures 3 and 4)}
We then introduce random variables as representations of uncertain outcomes:
\begin{itemize}[itemsep=-4pt, topsep=2pt]
    \item Discrete and continuous random variables.
    \item Probability mass and density functions, distribution and cumulative functions.
    \item Expected value, variance, covariance, and correlation.
    \item Joint and marginal distributions, independence of random variables.
    \item The normal and multivariate normal distributions and their importance in finance.
\end{itemize}
Illustrations include the modeling of asset returns, portfolio variance, and the concept of diversification through covariance matrices.

\subsection{Conditional Expectation (Lectures 4 and 5)}
Conditional expectation is introduced as a key tool for modeling information and forecasting:
\begin{itemize}[itemsep=-4pt, topsep=2pt]
    \item A quick definition of filtrations and sigma-algebra 
    \item Definition of conditional expectation and its interpretation as a projection.
    \item Definition of (sub/super) Martingale
\end{itemize}

\subsection{Convergence of Random Variables (Lectures 5 and 6)}
Finally, we discuss convergence results that justify the use of statistical estimation and simulation:
\begin{itemize}[itemsep=-4pt, topsep=2pt]
    \item Convergence in probability and almost sure convergence.
    \item Law of Large Numbers (LLN) and its role in empirical estimation.
    \item Central Limit Theorem (CLT)
    \item Monte Carlo simulation: generating random samples, estimating means, variances.
\end{itemize}
These concepts form the theoretical foundation for simulation-based methods used in modern quantitative finance.



\textbf{Applications:}
\begin{itemize}[itemsep=-4pt, topsep=2pt]
    \item Inferring the law of the underlying asset using Call options prices.
    \item Markov Chains: modeling credit rating transitions, regime-switching models, and stochastic state dynamics in risk processes.
    \item Monte Carlo Methods: simulating portfolio returns, estimating option prices.
    \item Law of Large Numbers and CLT in practice: convergence of simulated estimators, confidence intervals, and statistical validation of models.
    \item Linear Regression and Least Squares Estimation: parameter estimation, CAPM and multifactor models, calibration of linear pricing relationships.
\end{itemize}

\textbf{Suggested references:}
\begin{itemize}[itemsep=-4pt, topsep=2pt]
    \item Sheldon Ross, \textit{A First Course in Probability}.
    \item Robert Hogg \& Allen Craig, \textit{Introduction to Mathematical Statistics}.
    % \item Marek Musiela \& Marek Rutkowski, \textit{Martingale Methods in Financial Modelling} (for deeper reading).
\end{itemize}




\end{document}
